% !TeX root = ../papers/nc-p3.tex
\section{整な特異四次曲線を保つ極限}
\label{sec:integral-apolar-limits}
\begin{theorem}\label{thm:integral-apolar-alternatives}
楕円鉛筆の特殊化で、片側の定義二次式の極限が整な節点または尖点四次曲線$X$を定め、
特殊点が標準ASとする。可逆な二束を持つ定理\ref{thm:elliptic-pencil-integral-boundary}の表示か、
次の正規化上の表示のいずれかとなる。後者の必要なデータは$\nu:P=\mathbb P^1\to X$、
$A=H^0(X,\mathcal O_X(1))\subset H^0(P,\mathcal O(4))$、$B=H^0(P,\mathcal O(3))$と
$H^0(P,\mathcal O(7))$内の二平面$W$である。第一窓は積を$W$で割った写像の核の直交補で定まる。
後者の反対階数は8で、他側の反対核はねじれ三次曲線とその割線（節点）または接線（尖点）の合併を切る。
従って整な四次曲線を保つどの特殊化にも反対階数7の標準AS点はない。
後者は候補表示であり、定理\ref{thm:all-conductor-pencils-excluded}で全鉛筆を排除する。
\end{theorem}
\begin{proof}
特殊ファイバーまで延ばした完全交叉曲線族を$\mathcal X$、第一束を$L$とする。
付値環を必要なら完備化・有限拡大し、特殊点が非特異な切断$o$を選ぶ。
一般点上の第二束は$L_K(-o)(p)$と書ける。次数1の束は唯一の有効因子を持つからである。
固有性で$p$を切断に延ばし、そのイデアルを$J$とする。
$M=L(-o)\otimes\mathcal Hom(J,\mathcal O_{\mathcal X})$を用いる。
この双対は平坦で、ファイバー上で切断のイデアルの双対になる。
実際、双対列で$\mathcal O$の余核は$\mathcal Ext^1(\mathcal O_p,\mathcal O)$で、
切断上の可逆加群である。これは相対Gorenstein曲線における切断の双対性から従い、
付値環の剰余体への基底変換と可換である。局所的にも、切断を$(u,v)$、
曲線式を$u a+v b$で書けば、イデアルの双対は恒等写像と$(u,v)\mapsto(b,-a)$で生成される。
後者の類が切断上の自由な余核を生成し、剰余体上でも同じ表示を得る。

$p$が非特異点に落ちれば$M_0$は可逆である。特異点に落ちればその極大イデアルは導手であり、
逆イデアルは正規化環である（節点では二枝の値、尖点では一次係数の条件を外す）。
従って$M_0=\nu_*\mathcal O_P(3)$で、$L_0M_0=\nu_*\mathcal O_P(7)$となる。
いずれの場合も$h^0(M_0)=4,h^0(L_0M_0)=8$で$H^1=0$である。
可逆の場合の切断の積は既出の評価核により全射である。非可逆の場合は導手因子
$\Delta\subset P$（相異なる二点または二重点）で消える四次式全体が$A$に入る。
それに全三次式を掛ければ$\Delta$で消える七次式を全て得る。
$\Delta$上は単元となる$A$の元と三次式の積で任意の値を得るので、ここでも全射である。

命題\ref{prop:one-smooth-quartic-specialization}の格子の最大指数の議論を適用する。
低い指数の列に独立な$b,b'$があれば、その積$A b+A b'$が六次元以上であることを示す。
正規化上で共通因子を除いた二式の次数を$d\ge1$とする。積の核は
次数$4-d$の形式$h$で、互いに素な残りの両式を掛けるといずれも$A$に入るものからなる。
$d\ge3$なら次元は2以下である。$d=2$なら条件が恒等零になるには両式が$\Delta$で消えなければならない。
実際、$h$は$\Delta$上で任意の値を取り、$A$の貼り合わせ超平面は節点で両端の係数が非零、
尖点で微分の係数が非零である。どちらの局所長さ2の環にも、その超平面に入る非零イデアルはない。
互いに素なので少なくとも一条件がかかる。$d=1$なら二つの一次式が$\Delta$上の
全切断を張り、同じ理由で$h$自身が$\Delta$で消える。三次式のこの空間は二次元である。
よって主張が成立し、最大指数の列の評価の階数が2以下になる。

低い列が$s$一本なら、可逆な場合はその長さ4の零スキーム内の任意の長さ3の部分で消える
非零$x\in A$を取ればよい。長さ4のArtin環には長さ1のイデアルがあり、
長さ3への評価は四次元の$A$上で核を持つ。非可逆な場合は$s$が三次式で、
$s H^0(\mathcal O_P(1))$は四次式の超平面$A$と非零の$x$で交わる。
この$x$と全ての高い列の積は$s H^0(\mathcal O_P(4))/(sA)$の一次元に入る。
いずれも$x$の評価が低い一列を含めて階数2以下となり、反転窓の成長障害に反する。
従って最大・最小指数は等しく、特殊窓は実際に$A\otimes H^0(M_0)\to H^0(L_0M_0)/W$の表示を持つ。
可逆な場合、基点は純テンソルを与えるので元の整な鉛筆定理が適用できる。

非可逆な場合に反対切断を決める。$P$に別の直線$T$を$\Delta$で貼り合わせた種数1の曲線$Y$を使う。
二点で貼る場合は二成分の環、二重点の場合は接する二枝で、双対化層はいずれも自明である。
$T$を一点に縮めて$X$を得る射の引き戻しを$L'$とし、$M'$は$P$上で次数3、$T$上で次数1として貼る。
$\Delta$上の同定と$M'$の貼り合わせは任意に一つ選ぶ。
$T$上の一次切断は$\Delta$上の値で一意に決まるから、
$H^0(L')=A,H^0(M')=B,H^0(L'M')=H^0(\mathcal O_P(7))$で積は所与のものとなる。
両束は大域生成されるので、$W$に$Y$上の基点があればその評価が純テンソル関係になる。
従ってASの仮定のもとで$W$は$Y$上で基点自由である。

$|M'|$は$P$をねじれ三次曲線に、$T$を$\Delta$の張る直線に写す。
異なる二点の割線または二重点の接線と三次曲線の交叉はちょうど$\Delta$である。
その合併は二二次式完全交叉である。実際、三次曲線の三本の二次式の直線への制限は
$\Delta$で消える一次元空間を張り（標準単項式座標で確かめられる）、二本が残る。
可約二次式はねじれ三次曲線を含まないので共通因子はなく、合併が次数4を使い切る。
ゆえに両束の対称二乗は各々八次元の全二次切断に全射である。
$L'/M'$およびその逆は多重次数$(1,-1)$と$(-1,1)$で、次数1の成分だけでは
長さ2の貼り合わせに沿って消えられないため切断を持たない。

鉛筆完全列とSerre双対で切断の積のtraceを四線型形式にすると、
隣接対合は階数6で所与の六式を像とし、反対対合は階数8で両埋込みの二次核を持つ。
これは鉛筆完全列の計算であり、ここでは切断環のKoszul性を用いない。
特殊四窓が標準ASなら四連続の二次関係からの最上交叉は一次元なので、
この非零テンソルがその標準テンソルを張る。これで階数と反対核が従う。
\end{proof}

この非可逆な場合でもAS点の同型から曲線データを戻せる。整な側では特異点で消える二次切断、
他側では直線上で消える二次切断に制限する。両者は$P$上で固定因子を除けば
次数6と4の全切断であり、次元は7と5である。完全な八次元反対対合をここに制限すると階数は4以上である。
積は$T$上で零なので、正規化とねじれ三次曲線を同定するグラフへの制限を通る。
異なる二つの射影直線同型のグラフは長さ2で交わり、合併のイデアルで$(6,4)$をひねると
$(4,2)$で$H^1=0$となる。従って二つの線型包の共通部分の行列階数は高々2である。
これで同定が一意に決まり、$W$も隣接核から戻る。逆にこの正規化上の順序付きデータの同型は六式を保つ。

\paragraph{切断環の三次障害。}
非可逆模型の$Y$上の交互切断環は二次生成関係だけでは表示できない。
節点では$A=\langle1+t^4,t,t^2,t^3\rangle$、尖点では
$A=\langle1,t^2,t^3,t^4\rangle$に揃えられる。前者では二枝の座標のスカラーを、
後者では$t/(1+ct)$の四次形式への作用を用いて貼り合わせ係数を規格化する。
$B=\langle1,t,t^2,t^3\rangle$として、$B,A,B$の二次乗法核から作る三次商は両方で13次元、
実際の切断は12次元である。三次の積は全射である。主成分では$AB$が全七次式で
さらに$B$を掛けると全十次式となる。主成分上で消えて$T$上で$\Delta$に沿って消える
一つの二次式は、ねじれ三次曲線上で消え直線上では消えない二次式を取り、
$T$で非零の$A$の切断を掛ければ得られる。
$A,B,A$の二次関係からの三次商は12次元となる。
\path{computations/verify_conductor_pencil.py}はそれぞれの二次乗法核を64行の行列に入れてこの次元を
有理数上で検証し、各型の基点自由な一鉛筆で隣接・反対階数6,8と両三次交叉の次元4を検査する。
曲線の交互切断環の非二次性から、二次の鉛筆を割った関係列の非AS性は導けない。

\begin{proposition}[非可逆模型の全鉛筆における有限双対]\label{prop:conductor-pencil-frobenius}
上の正規化表示で、幾何的な二成分曲線$Y$上で鉛筆$W$が基点自由とする。
全てのこのような$W$で、二次関係列の二次双対は次元列$1,4,6,4,1$の有限Frobenius代数である。
第三切断階数は8で、順序付き正規化データによる窓の同型判定も成立する。
\end{proposition}
\begin{proof}
まず二次の積だけで定めた交互切断環の二次表示を考える。
$A,B,A$の三次商は上の計算で12次元で、実際の切断への写像は全射である。
実際$AB$は全七次式であり、$A$と七次式の積も導手部分と$\Delta$上の値を分ければ全十一
次式を張る。この位置では$T$への制限は次数1で、主成分での値から決まる。
従って三次商と切断が一致し、$W$を両辺から割れば$W A$の八次元が消えて四次元になる。
$W A$の次元は、基点自由鉛筆の列で$A$に対応する束へ積を取り、積の核が
$H^0((M')^{-1})=0$となるためである。$B$についても同様に$\dim W B=8$である。

$B,A,B$の三次商から切断への全射の核は一次元である。その余分な方向は、
切断の積と合わせて、次の線型汎関数で検出できる。
\[
 b\otimes a\otimes b'\longmapsto
 a|_T\,\det(b|_T,b'|_T).
\]
ここで$L'$は$T$上で自明であり、$M'|_T=\mathcal O(1)$の二切断の行列式を取る。
この式は左の二次積の後にも右の二次積の後にも定義できるので、二次関係を消す。
しかも外側の二文字の交換で符号が変わり、通常の切断の積では得られない非零の汎関数である。
$w\in W$を左右の二次積へ互いに反転する形で持ち上げれば、
$w b-b w$のこの値は$2\det(w|_T,b|_T)$である。
基点自由性により$W$の$T$への制限は二次元であり、行列式を非零にできる。
従って$W$を両辺から割る際には、切断上で$W B$が消えるのに加えてこの一次元も消える。
結果は$13-8-1=4$次元となる。

上のtraceテンソルはどの三次交叉でもその全体を一文字縮約で張る。
縮約が四次元となるのは、文字が三次元以下で足りるならその側の対称二乗を通る反対階数が
6以下となるためである。四次交叉の元は、最初の三文字または最後の三文字をまとめると、
traceテンソルに第一位置の自己準同型を施したものでも、第四位置の自己準同型を施したものでもある。
後者は第一・第三位置で対称なので、前者の第一位置の作用も八次元の対称形式で自己随伴になる。
定理\ref{thm:elliptic-pencil-reconstruction}の次元7の上界の証明は非可逆な自己準同型にも通る。
従って第一位置の作用はスカラーで、四次交叉は一次元である。
五次交叉はこの四次テンソルを左・右に一文字延ばした二空間の共通部分である。
非零の共通元は第一位置だけを切ると、左が階数4、右が階数1となりえない。
従って五次以上は零である。二次双対はこれでtraceから得る完全対合の商と全次数で一致する。
両反対核とグラフの復元にはこのtraceとその階数だけを使ったので、順序付き同型判定も従う。
\end{proof}

純テンソル関係が無い条件もこの模型では幾何的に読める。
$H^0(\mathcal O_P(7))^*$内の汎関数が$A\otimes B$で階数1なら、
ある超平面$A'\subset A$について$A' B$を消す。
$A'$が$P$上で基点自由なら一般の二切断が互いに素な四次式で、
三次式との積は八次元全体を張るので矛盾する。従って$A'$は点での評価核である。
特異点の原像以外では、$|A|$はそこで埋込みで他の点と分離するので共通因子はその一点だけである。
一般二切断から共通因子を除いた互いに素な三次式と$B$の積は全六次式を張り、
元の汎関数はその点の評価に限られる。特異点の原像では$A'$は$\Delta$で消える全四次式なので、
汎関数は$\Delta$への制限を通る二次元空間で尽くされる。
この空間の全ての非零元は第一因子が特異点での評価の倍数で、実際に階数1である。
ゆえに射影化した階数1の軌跡は七次有理正規曲線と$\Delta$の張る直線の合併で、
$W^\perp$がこれを避けることは$Y$上の基点自由性と同値である。
また一つの文字の評価核が二次元以上なら、その文字と相手側の切断の積の中に$W$全体が入る。
主成分上では一文字が必ず零点を持つので、これも基点自由性に反する。

\begin{corollary}[七次で排除される二つの鉛筆]\label{cor:conductor-degree-seven}
節点型と尖点型の双方で、$W=\langle1+2t+4t^2+7t^6+3t^7,\ t+4t^4+t^7\rangle$
は基点自由で有限双対がFrobeniusであるが、標準ASではない。
\end{corollary}
\begin{proof}
両七次式の終結式は非零で、例外直線上でも二切断は独立である。
上記コードは、七次式の単項式座標における積の係数行列に$W^\perp$の行基底を掛け、
その六行と転置順の六行を交互の関係として使う。
既存の\path{computations/two_window_ext.py}の経路商で、次数$d-2$の基底に
最後の二文字の六関係を掛けた全ての像を、次数$d-1$の基底に一文字掛けた空間から消去する。
有理数上の両始点の次元は六次まで$1,4,10,20,35,56,84$であるが、
七次では$A$始点で120、$B$始点で121となる。標準値$\binom{10}{3}=120$に反する。
\end{proof}

節点型ではさらに$W_u=\langle1,t^7+u t^3\rangle$の全ての$u\in k$で標準AS性を排除できる。
主成分と例外直線で基点がない。上記コードは$1=0,t^7=-u t^3$で七次式の積を商に写し、
有理関数体$\mathbb Q(u)$上で$B$始点の七次次元121を得る。特殊化で一定階数の六関係が定める
各経路商の次元は上半連続なので、全ての有限パラメータで七次元数は121以上である。
